
\documentclass{article}
\usepackage{colortbl}
\usepackage{makecell}
\usepackage{multirow}
\usepackage{supertabular}

\begin{document}

\newcounter{utterance}

\centering \large Interaction Transcript for game `cladder', experiment `full\_v1.5\_default', episode 2677 with Qwen3.5{-}27B.
\vspace{24pt}

{ \footnotesize  \setcounter{utterance}{1}
\setlength{\tabcolsep}{0pt}
\begin{supertabular}{c@{$\;$}|p{.15\linewidth}@{}p{.15\linewidth}p{.15\linewidth}p{.15\linewidth}p{.15\linewidth}p{.15\linewidth}}
    \# & \multicolumn{2}{c}{Player} && \multicolumn{2}{c}{Game Master} \\
    \hline

    \theutterance \stepcounter{utterance}  
    & & & \multicolumn{4}{p{0.6\linewidth}}{
        \cellcolor[rgb]{0.9,0.9,0.9}{
            \makecell[{{p{\linewidth}}}]{
                \texttt{\tiny{[P1$\langle$GM]}}
                \texttt{You are an expert in causal inference. The following question is not a typical commonsense query, but rather a meticulously designed question created by a professor specializing in causal inference, intended to assess the students' mastery of the course content.} \\
\texttt{Imagine a self{-}contained, hypothetical world with only the following conditions, and without any unmentioned factors or causal relationships: Gene has a direct effect on smoking and lung cancer. Pollution has a direct effect on smoking and lung cancer. Smoking has a direct effect on lung cancer. Pollution is unobserved. Method 1: We look at how gene correlates with lung cancer case by case according to smoking. Method 2: We look directly at how gene correlates with lung cancer in general. To understand how gene affects lung cancer, is it more correct to use the Method 1 than Method 2?} \\
\texttt{Start your answer with 'Yes' or 'No', followed by additional reasoning or evidence to support your explanation.} \\
            }
        }
    }
     \\ \\

    \theutterance \stepcounter{utterance}  
    & \multicolumn{4}{p{0.6\linewidth}}{
        \cellcolor[rgb]{0.561,0.737,0.561}{
            \makecell[{{p{\linewidth}}}]{
                \texttt{\tiny{[P1$\rangle$GM]}}
                \texttt{No} \\
\\ 
\texttt{In this hypothetical scenario, **Method 2** (looking at the direct correlation between Gene and Lung Cancer in general) is the correct approach to estimate the total causal effect of Gene on Lung Cancer, whereas **Method 1** (stratifying by Smoking) introduces bias.} \\
\\ 
\texttt{Here is the causal reasoning based on the structure provided:} \\
\\ 
\texttt{1.  **Causal Structure**:} \\
\texttt{    *   \$Gene \textbackslash{}rightarrow Smoking\$} \\
\texttt{    *   \$Gene \textbackslash{}rightarrow LungCancer\$} \\
\texttt{    *   \$Pollution \textbackslash{}rightarrow Smoking\$} \\
\texttt{    *   \$Pollution \textbackslash{}rightarrow LungCancer\$} \\
\texttt{    *   \$Smoking \textbackslash{}rightarrow LungCancer\$} \\
\texttt{    *   \$Pollution\$ is unobserved.} \\
\\ 
\texttt{2.  **Analysis of Method 2 (Marginal Association)**:} \\
\texttt{    *   This method estimates the association \$P(LungCancer | Gene)\$.} \\
\texttt{    *   The paths from \$Gene\$ to \$LungCancer\$ are:} \\
\texttt{        *   Direct path: \$Gene \textbackslash{}rightarrow LungCancer\$ (Causal)} \\
\texttt{        *   Indirect path: \$Gene \textbackslash{}rightarrow Smoking \textbackslash{}rightarrow LungCancer\$ (Causal)} \\
\texttt{    *   Is there a confounding path? A confounding path would require a common cause of \$Gene\$ and \$LungCancer\$. The only common cause mentioned is \$Pollution\$, but \$Pollution\$ does **not** affect \$Gene\$. Therefore, there is no backdoor path from \$Gene\$ to \$LungCancer\$.} \\
\texttt{    *   Consequently, the marginal association in Method 2 captures the **total causal effect** of Gene on Lung Cancer (both direct and indirect via smoking) without bias.} \\
\\ 
\texttt{3.  **Analysis of Method 1 (Stratification by Smoking)**:} \\
\texttt{    *   This method attempts to estimate the effect of \$Gene\$ on \$LungCancer\$ while holding \$Smoking\$ constant (e.g., \$P(LungCancer | Gene, Smoking)\$).} \\
\texttt{    *   \$Smoking\$ is a **mediator** on the causal path from \$Gene\$ to \$LungCancer\$.} \\
\texttt{    *   Conditioning on a mediator blocks the indirect causal path (\$Gene \textbackslash{}rightarrow Smoking \textbackslash{}rightarrow LungCancer\$). If the goal is to understand "how Gene affects Lung Cancer" (implying the total effect), blocking this path yields an incomplete estimate (the direct effect only).} \\
\texttt{    *   **Crucially**, \$Smoking\$ is also a **collider** on the path \$Gene \textbackslash{}rightarrow Smoking \textbackslash{}leftarrow Pollution \textbackslash{}rightarrow LungCancer\$.} \\
\texttt{        *   Normally, the path \$Gene \textbackslash{}leftarrow \textbackslash{}dots \textbackslash{}rightarrow Pollution \textbackslash{}rightarrow LungCancer\$ is blocked because \$Gene\$ and \$Pollution\$ are independent.} \\
\texttt{        *   However, by conditioning on \$Smoking\$ (the collider),} \\
            }
        }
    }
    & & \\ \\

    \theutterance \stepcounter{utterance}  
    & & & \multicolumn{2}{p{0.3\linewidth}}{
        \cellcolor[rgb]{0.9,0.9,0.9}{
            \makecell[{{p{\linewidth}}}]{
                \texttt{\tiny{[GM$|$GM]}}
                \texttt{No} \\
            }
        }
    }
    & & \\ \\

    \theutterance \stepcounter{utterance}  
    & & & \multicolumn{2}{p{0.3\linewidth}}{
        \cellcolor[rgb]{0.9,0.9,0.9}{
            \makecell[{{p{\linewidth}}}]{
                \texttt{\tiny{[GM$|$GM]}}
                \texttt{no} \\
            }
        }
    }
    & & \\ \\

    \theutterance \stepcounter{utterance}  
    & & & \multicolumn{2}{p{0.3\linewidth}}{
        \cellcolor[rgb]{0.9,0.9,0.9}{
            \makecell[{{p{\linewidth}}}]{
                \texttt{\tiny{[GM$|$GM]}}
                \texttt{game\_result = WIN} \\
            }
        }
    }
    & & \\ \\

\end{supertabular}
}

\end{document}
